% M1_03: Section 2.3, justification of Laplace noise (technical role and economic sources).
% REPLACES: main.tex lines 337-340 ("Here \(V_T\) is the terminal payoff ..." to "... also considers logistic noise.")
% The replacement keeps the definitions of \(V_T\) and \(Z\) and the pointer to logistic noise.
% It refers to the class result by its label, result-proposition-a-5.
% NUM placeholders: none.
% New \bibitem requested: ref-SmithSorensen2000, ref-Clark1973, ref-Cont2001 (keys in BRIEF Section 7).
% New symbols: none. \(m(\pi)\) and \(M(\pi)\) use the notation sheet of the brief.

Here \(V_T\) is the terminal payoff of a target share and \(Z\) is
independent noise demand with a Laplace density. The Laplace density
has one technical role. Up to a constant, its logarithm is
\(-|z|/b\), with slope \(\pm 1/b\). For each pair of feasible orders,
the likelihood ratio of a flow then lies in \([e^{-2/b},e^{2/b}]\).
No order flow is conclusive, in the sense of bounded private beliefs
(\citeproc{ref-SmithSorensen2000}{Smith and S{\o}rensen 2000}). Each
posterior stays in \([m(\pi),M(\pi)]\), so the investor keeps a
residual advantage at each flow. A change in the order also changes
the flow density by at most the fraction \(1/b\) of its level. With
\(b>1\), this score bound gives a lower bound on marginal trading
profit on the full order interval. The bound covers mixed orders and
each deviation, and it gives uniqueness in
\hyperref[result-proposition-2]{Proposition 2}. Under full orders the
posterior log odds equal
\(\operatorname{logit}\pi+(2/b)\max\{-1,\min\{x,1\}\}\). Between the
two orders the log odds are linear in flow, as under Gaussian noise
with variance \(b\). Outside the two orders the log odds are flat.
The linear part gives closed forms. The flat part keeps the market
uncertain after extreme flows. Gaussian noise keeps the linear log
odds and loses the bounds: an extreme flow all but reveals the class,
the residual advantage vanishes there, and the global argument fails.
Each density with \(|(\log f)'|\le L<1\) has the same two properties.
\hyperref[result-proposition-a-5]{Proposition A.5} states a version of
\hyperref[result-proposition-2]{Proposition 2} for this bounded-score
class. The logistic density is one member, and
\hyperref[sec-extensions]{Section 5} reports that case. Within the
class at a fixed bound, Laplace noise is the extreme member: it is
the least dispersed law, and at the same orders it gives the most
informative price.

Three standard models give a Laplace law. Net noise demand is Laplace
if noise buy volume and noise sell volume are independent exponential
variables with equal means. It is Laplace if it is Gaussian with an
exponential variance, as with a random number of liquidity traders.
It is Laplace if noise flow follows a Brownian motion up to an
independent exponential time, such as the arrival of public news.
The last model is the subordination device of Clark
(\citeproc{ref-Clark1973}{1973}). Each member of the class has
tails at least as fat as an exponential tail. Asset returns also have
tails fatter than Gaussian tails
(\citeproc{ref-Cont2001}{Cont 2001}). The class claim is therefore
about tail weight and bounded inference, not about a fit of one law
to noise demand.
